% GENERATED FILE. Edit canonical lessons, then run npm run generate:books.
% canonical-source-hash: fnv1a64:e8d37b4263588210
% canonical-lessons: GE-C402-tanken, GE-C402-volltanken, GE-C402-benzin, GE-C402-zapfsaeule, GE-C402-oelwechsel

\chapter{An der Tankstelle}
\label{ch:ge-an-der-tankstelle}
\hlchaptermodality{402}

\begin{chapteropening}
\textbf{By the end of this chapter:} \emph{I can fill up the car: filling up, filling the tank, petrol, the pump and the oil change.}

Complete the last lesson of chapter 402: I can fill up the car: filling up, filling the tank, petrol, the pump and the oil change.
\end{chapteropening}

\section[tanken]{tanken — to fill up (with fuel), to get petrol}

\label{lesson:GE-C402-tanken}



\begin{quote}
Before the new one: say the German for to buckle up, then the German for the front-seat passenger.
\end{quote}

\subsection*{You'll want to know: tanken}
\textbf{tanken} — \textquotedblleft{}to fill up (with fuel), to get petrol\textquotedblright{}.

\begin{quote}
Auf dem Weg ans Meer tanken wir einmal.
\end{quote}

\begin{quote}
\emph{On the way to the sea we fill up once.}
\end{quote}

\emph{Er tankt}, \emph{er hat getankt}. \emph{Die Tankstelle} is the petrol station.

\begin{cousinweb}
\emph{tanken} is built on \emph{Tank} 'tank', which was borrowed from English \emph{tank} (EtymDB). EtymDB traces the English word further back to Portuguese \emph{tanque} 'water tank'.
\end{cousinweb}

\begin{grammarlens}[title={What you've built}]
One more word for telling, explaining and describing at B1.
\end{grammarlens}

\subsection*{Your turn}
\begin{itemize}
\raggedright
  \item \emph{Say it:} \emph{tanken}
  \item \emph{Say it:} \emph{tanken}, once more
  \item \emph{Say it:} a sentence of your own with \emph{tanken}
  \item \emph{Recall:} say the German for to buckle up, then the German for the front-seat passenger, then say \emph{tanken} again
\end{itemize}

\begin{tcolorbox}[breakable,colback=teal!4,colframe=teal!35!black,title={Before you move on}]
What does tanken mean? (\textquotedblleft{}To fill up (with fuel), to get petrol\textquotedblright{}.) Say it once more.
\end{tcolorbox}
\section[volltanken]{volltanken — to fill the tank right up}

\label{lesson:GE-C402-volltanken}



\begin{quote}
Before the new one: say the German for the front-seat passenger, then the German for to fill up.
\end{quote}

\subsection*{You'll want to know: volltanken}
\textbf{volltanken} — \textquotedblleft{}to fill the tank right up\textquotedblright{}.

\begin{quote}
Vor der langen Reise tankt er das Auto voll.
\end{quote}

\begin{quote}
\emph{Before the long trip he fills the car right up.}
\end{quote}

Separable: \emph{ich tanke voll}, \emph{ich habe vollgetankt}.

\begin{cousinweb}
\emph{volltanken} is built from \emph{voll} 'full', an inherited Germanic word, and \emph{tanken}, from \emph{Tank}, which was borrowed from English \emph{tank} (EtymDB).
\end{cousinweb}

\begin{grammarlens}[title={What you've built}]
One more word for telling, explaining and describing at B1.
\end{grammarlens}

\subsection*{Your turn}
\begin{itemize}
\raggedright
  \item \emph{Say it:} \emph{volltanken}
  \item \emph{Say it:} \emph{volltanken}, once more
  \item \emph{Say it:} a sentence of your own with \emph{volltanken}
  \item \emph{Recall:} say the German for the front-seat passenger, then the German for to fill up, then say \emph{volltanken} again
\end{itemize}

\begin{tcolorbox}[breakable,colback=teal!4,colframe=teal!35!black,title={Before you move on}]
What does volltanken mean? (\textquotedblleft{}To fill the tank right up\textquotedblright{}.) Say it once more.
\end{tcolorbox}
\section[das Benzin]{das Benzin — petrol (US: gas)}

\label{lesson:GE-C402-benzin}



\begin{quote}
Before the new one: say the German for to fill up, then the German for to fill the tank right up.
\end{quote}

\subsection*{You'll want to know: das Benzin}
\textbf{das Benzin} — \textquotedblleft{}petrol (US: gas)\textquotedblright{}.

\begin{quote}
In den Ferien ist das Benzin wieder teurer.
\end{quote}

\begin{quote}
\emph{In the holidays petrol is more expensive again.}
\end{quote}

No plural in everyday use. \emph{Bleifrei} — unleaded.

\begin{cousinweb}
\emph{Benzin} is built on \emph{Benzoe} 'benzoin', a sweet-smelling resin, with the chemist's ending \emph{-in} (EtymDB). The name of the resin is usually traced to Arabic.
\end{cousinweb}

\begin{grammarlens}[title={What you've built}]
One more word for telling, explaining and describing at B1.
\end{grammarlens}

\subsection*{Your turn}
\begin{itemize}
\raggedright
  \item \emph{Say it:} \emph{das Benzin}
  \item \emph{Say it:} \emph{das Benzin}, once more
  \item \emph{Say it:} a sentence of your own with \emph{das Benzin}
  \item \emph{Recall:} say the German for to fill up, then the German for to fill the tank right up, then say \emph{das Benzin} again
\end{itemize}

\begin{tcolorbox}[breakable,colback=teal!4,colframe=teal!35!black,title={Before you move on}]
What does das Benzin mean? (\textquotedblleft{}Petrol (US: gas)\textquotedblright{}.) Say it once more.
\end{tcolorbox}
\section[die Zapfsäule]{die Zapfsäule — the petrol pump (at a filling station)}

\label{lesson:GE-C402-zapfsaeule}



\begin{quote}
Before the new one: say the German for to fill the tank right up, then the German for petrol.
\end{quote}

\subsection*{You'll want to know: die Zapfsäule}
\textbf{die Zapfsäule} — \textquotedblleft{}the petrol pump (at a filling station)\textquotedblright{}.

\begin{quote}
Wir halten an der Zapfsäule neben dem kleinen Laden.
\end{quote}

\begin{quote}
\emph{We stop at the pump next to the little shop.}
\end{quote}

Plural: \emph{die Zapfsäulen}. \emph{Zapfen} means to draw off a liquid, like beer from a barrel.

\begin{cousinweb}
\emph{Zapfsäule} is built from \emph{zapfen} 'to draw off a liquid', from \emph{Zapfen} 'plug, bung', and \emph{Säule} 'pillar', an inherited Germanic word (EtymDB lists \emph{Säule}).
\end{cousinweb}

\begin{grammarlens}[title={What you've built}]
One more word for telling, explaining and describing at B1.
\end{grammarlens}

\subsection*{Your turn}
\begin{itemize}
\raggedright
  \item \emph{Say it:} \emph{die Zapfsäule}
  \item \emph{Say it:} \emph{die Zapfsäule}, once more
  \item \emph{Say it:} a sentence of your own with \emph{die Zapfsäule}
  \item \emph{Recall:} say the German for to fill the tank right up, then the German for petrol, then say \emph{die Zapfsäule} again
\end{itemize}

\begin{tcolorbox}[breakable,colback=teal!4,colframe=teal!35!black,title={Before you move on}]
What does die Zapfsäule mean? (\textquotedblleft{}The petrol pump (at a filling station)\textquotedblright{}.) Say it once more.
\end{tcolorbox}
\section[der Ölwechsel]{der Ölwechsel — the oil change}

\label{lesson:GE-C402-oelwechsel}



\begin{quote}
Before the new one: say the German for petrol, then the German for the petrol pump.
\end{quote}

\subsection*{You'll want to know: der Ölwechsel}
\textbf{der Ölwechsel} — \textquotedblleft{}the oil change\textquotedblright{}.

\begin{quote}
Nach dem Ölwechsel fährt das Auto wieder ruhig.
\end{quote}

\begin{quote}
\emph{After the oil change the car runs smoothly again.}
\end{quote}

Same form in the plural.

\begin{cousinweb}
\emph{Ölwechsel} is built from \emph{Öl} 'oil', an early borrowing through Old High German \emph{oli} from Latin \emph{oleum} (EtymDB), and \emph{Wechsel} 'change', inherited from Old High German \emph{wehsal} (EtymDB).
\end{cousinweb}

\begin{grammarlens}[title={What you've built}]
One more word for telling, explaining and describing at B1.
\end{grammarlens}

\subsection*{Your turn}
\begin{itemize}
\raggedright
  \item \emph{Say it:} \emph{der Ölwechsel}
  \item \emph{Say it:} \emph{der Ölwechsel}, once more
  \item \emph{Say it:} a sentence of your own with \emph{der Ölwechsel}
  \item \emph{Recall:} say the German for petrol, then the German for the petrol pump, then say \emph{der Ölwechsel} again
\end{itemize}

\begin{tcolorbox}[breakable,colback=teal!4,colframe=teal!35!black,title={Before you move on}]
What does der Ölwechsel mean? (\textquotedblleft{}The oil change\textquotedblright{}.) Say it once more.
\end{tcolorbox}
